%============================================================
% CH1 · PATTERN 3 — EXERCISES, HINTS, ANSWERS
%============================================================
\subsection{Pattern 3: Exercises and Going Further}
\label{ssec:ch1-p3-ex}

\begin{patternbox}{Pattern 3 Summary Card}
\begin{tabularx}{\linewidth}{@{}>{\bfseries}l X@{}}
Recognise & reverse, rotate, cyclic shift; two blocks swap places ($AB \to BA$); ``is it a rotation of \ldots?'' \\
Core loops & reverse $[l,r)$: swap $a[l], a[r-1]$ and move inwards; rotate right by $k$: $k \bmod n$, reverse all, reverse $[0,k)$, reverse $[k,n)$ \\
Tools & cycles of $i \to (i+k) \bmod n$ ($\gcd(n,k)$ of them); a start offset instead of moving data; rotations $=$ windows of $aa$; at most one circular drop $\iff$ rotated sorted \\
Cost & $\bigO{n}$ time, $\bigO{1}$ extra memory \\
Top traps & no $k \bmod n$; left/right confused; in-place moves that overwrite unread values; forgetting the seam pair $(a[n-1], a[0])$; negative \texttt{\%} \\
\end{tabularx}
\end{patternbox}

\exlegend
\begin{enumerate}[leftmargin=*, label=\textbf{E3.\arabic*}, itemsep=6pt]
    \item \exeasy\ \pconcept \exprob{Reverse String II}{LeetCode 541}{https://leetcode.com/problems/reverse-string-ii/}\ \lc{Easy}. Split a string of length $n \le 10^4$ into consecutive blocks of $2k$ characters (the last block may be shorter). In every block reverse the first $k$ characters (all of them if fewer than $k$ remain) and keep the rest.
    \item \exeasy\ \precent \exprob{Matrix Similarity After Cyclic Shifts}{LeetCode 2946}{https://leetcode.com/problems/matrix-similarity-after-cyclic-shifts/}\ \lc{Easy}. An $m \times n$ matrix ($m, n \le 25$) undergoes $k \le 50$ steps; in each step the even-indexed rows are cyclically shifted one place to the left and the odd-indexed rows one place to the right. Is the final matrix equal to the original one?
    \item \exeasy\ \pconcept \cf{900}\ \exprob{Mainak and Array}{Codeforces 1726A}{https://codeforces.com/problemset/problem/1726/A}. You must perform exactly one operation: choose a subarray $[l, r]$ and rotate it cyclically by any amount. Maximise $a_n - a_1$. ($n \le 2000$ per test.)
    \item \exmedium\ \pclassic \exprob{Rotate Function}{LeetCode 396}{https://leetcode.com/problems/rotate-function/}\ \lc{Medium}. For the array $B_k$ obtained by rotating $a$ to the right by $k$, let $F(k) = \sum_i i \cdot B_k[i]$. Return $\max_{0 \le k < n} F(k)$. $n \le 10^5$, $|a_i| \le 100$.
    \item \exeasy\ \ptricky \exprob{Make Two Arrays Equal by Reversing Subarrays}{LeetCode 1460}{https://leetcode.com/problems/make-two-arrays-equal-by-reversing-subarrays/}\ \lc{Easy}. In one step you may reverse any non-empty subarray of \texttt{arr}. Can \texttt{arr} be made equal to \texttt{target} (same length $n \le 1000$, values $\le 1000$)?
    \item \exmedium\ \ptricky \cf{1300}\ \exprob{Sort the Array}{Codeforces 451B}{https://codeforces.com/problemset/problem/451/B}. Given $n \le 10^5$ \emph{distinct} integers, decide whether reversing exactly one segment $[l, r]$ makes the array increasing; if yes, print such a segment ($l = r$ is allowed).
\end{enumerate}

\subsubsection*{Hints}
\begin{enumerate}[leftmargin=*, label=\textbf{E3.\arabic*}]
    \item The blocks start at $i = 0, 2k, 4k, \dots$; which range is reversed in the block starting at $i$? Watch the end of the string.
    \item Shifting a row of length $n$ left $k$ times is a left rotation by $k \bmod n$. When is a row equal to its own rotation? Does the direction matter for that question?
    \item Only $a_1$ and $a_n$ matter. Which values can a single rotation bring to the last position while $a_1$ stays? To the first while $a_n$ stays? What if both ends move?
    \item Compute $F(0)$ directly. Going from $B_{k-1}$ to $B_k$, every element moves one index to the right except one, which wraps from index $n-1$ to $0$.
    \item What can two reversals of length 2 do? Which arrangements can then be reached?
    \item The segment to reverse must start at the first place where the order breaks. How far must it extend? After reversing, one check is enough.
\end{enumerate}

\subsubsection*{Answers}

\begin{solutionbox}[E3.1 Reverse String II]
For every block start $i = 0, 2k, 4k, \dots$ reverse the half-open range $[i, \min(i+k, n))$ with the two-index loop (here \texttt{std::reverse}). The \texttt{min} handles both short endings at once: fewer than $k$ characters left (reverse them all) and between $k$ and $2k$ left (reverse the first $k$).
\begin{lstlisting}[style=cpp]
#include <bits/stdc++.h>
using namespace std;

class Solution {
public:
    string reverseStr(string s, int k) {
        int n = s.size();
        for (int i = 0; i < n; i += 2 * k)
            reverse(s.begin() + i, s.begin() + min(i + k, n));   // [i, min(i+k, n))
        return s;
    }
};

int main() {
    cout << Solution().reverseStr("abcdefg", 2) << " "
         << Solution().reverseStr("abcd", 2) << "\n";
    // bacdfeg bacd
}
\end{lstlisting}
\begin{center}
\begin{tikzpicture}[x=1cm, y=1cm]
    \dparray{s}{0.9}{before}{a/aread, b/aread, c/acell, d/acell, e/aread, f/aread, g/acell}
    \dparray{t}{-0.4}{after}{b/awrite, a/awrite, c/acell, d/acell, f/awrite, e/awrite, g/acell}
    \arange{t0}{t3}{block 1}{memoorange}
    \arange{t4}{t6}{block 2}{memoorange}
\end{tikzpicture}
\end{center}
Total work $\bigO{n}$. Both official samples reproduced.
\end{solutionbox}

\begin{solutionbox}[E3.2 Matrix Similarity After Cyclic Shifts]
A row is shifted one place per step, so after $k$ steps an even row is rotated left by $k$ and an odd row right by $k$ --- no need to simulate. A row $r$ equals its left rotation by $k$ iff $r[j] = r[(j + k) \bmod n]$ for every $j$; and that holds iff it equals its \emph{right} rotation by $k$ (substitute $j \to j - k$). So the direction is a red herring, and the trap is to simulate it carefully when the answer does not depend on it.
\begin{lstlisting}[style=cpp]
#include <bits/stdc++.h>
using namespace std;

class Solution {
public:
    bool areSimilar(vector<vector<int>>& mat, int k) {
        for (auto& row : mat) {
            int n = row.size();
            for (int j = 0; j < n; j++)
                if (row[j] != row[(j + k) % n]) return false;   // rotation as an index shift
        }
        return true;
    }
};

int main() {
    vector<vector<int>> m1 = {{1, 2, 3}, {4, 5, 6}, {7, 8, 9}};
    vector<vector<int>> m2 = {{1, 2, 1, 2}, {5, 5, 5, 5}, {6, 3, 6, 3}};
    cout << Solution().areSimilar(m1, 4) << Solution().areSimilar(m2, 2) << "\n";   // 01
}
\end{lstlisting}
$\bigO{mn}$ time, $\bigO{1}$ memory. Official samples $\to$ false, true, true; checked against a literal simulation of the $k$ steps on random $0/1$ matrices.
\end{solutionbox}

\begin{solutionbox}[E3.3 Mainak and Array]
Three kinds of rotation can help, and every optimal operation is one of them:
\begin{itemize}[leftmargin=*]
    \item keep $a_1$, rotate a suffix so that its maximum lands at position $n$: value $\max_{i \ge 2} a_i - a_1$;
    \item keep $a_n$, rotate a prefix so that its minimum lands at position 1: value $a_n - \min_{i \le n-1} a_i$;
    \item rotate the whole array (both ends change): the new ends are neighbours $a_{i-1}$ (last) and $a_i$ (first) of the old array: value $\max_i (a_{i-1} - a_i)$, including the rotation by 0, i.e.\ $a_n - a_1$.
\end{itemize}
A subarray that contains neither end changes nothing; one that contains exactly one end is covered by the first two cases; one that contains both is a rotation of everything.
\begin{lstlisting}[style=cpp]
#include <bits/stdc++.h>
using namespace std;

int main() {
    int t;
    cin >> t;
    while (t--) {
        int n;
        cin >> n;
        vector<int> a(n);
        for (auto& x : a) cin >> x;
        int best = a[n - 1] - a[0];                                    // rotate by 0
        for (int i = 1; i < n; i++) best = max(best, a[i] - a[0]);     // a_1 fixed
        for (int i = 0; i < n - 1; i++) best = max(best, a[n - 1] - a[i]);  // a_n fixed
        for (int i = 1; i < n; i++) best = max(best, a[i - 1] - a[i]); // whole array
        cout << best << "\n";
    }
}
\end{lstlisting}
Sample $[1,3,9,11,5,7]$: the three cases give $11-1 = 10$, $7-1 = 6$ and $\max(a_{i-1}-a_i) = 11-5 = 6$, so 10. Sample $[2,1,8,1]$: $6$, $0$ and $8-1 = 7$ $\Rightarrow 7$. All five sample answers $10, 0, 990, 7, 4$ reproduced, and checked against trying every subarray and every shift for $n \le 6$.
\end{solutionbox}

\begin{solutionbox}[E3.4 Rotate Function]
Let $S = \sum a_i$. From $B_{k-1}$ to $B_k$ every element moves one index to the right, adding $S$ in total --- except the element that was at index $n-1$, namely $a[n-k]$, which wraps to index 0: it contributed $(n-1)a[n-k]$ and should have gained one more $a[n-k]$, but now contributes $0$. Hence
\[
F(k) = F(k-1) + S - n \cdot a[n-k], \qquad k = 1, \dots, n-1 .
\]
\begin{lstlisting}[style=cpp]
#include <bits/stdc++.h>
using namespace std;
typedef long long ll;

class Solution {
public:
    int maxRotateFunction(vector<int>& nums) {
        int n = nums.size();
        ll S = 0, F = 0;
        for (int i = 0; i < n; i++) { S += nums[i]; F += (ll)i * nums[i]; }   // F(0)
        ll best = F;
        for (int k = 1; k < n; k++) {
            F += S - (ll)n * nums[n - k];        // a[n-k] wraps from the end to index 0
            best = max(best, F);
        }
        return best;
    }
};

int main() {
    vector<int> a = {4, 3, 2, 6};
    cout << Solution().maxRotateFunction(a) << "\n";   // 26
}
\end{lstlisting}
Sample $[4,3,2,6]$: $S = 15$, $F(0) = 25$, $F(1) = 25 + 15 - 4\cdot 6 = 16$, $F(2) = 16 + 15 - 4\cdot 2 = 23$, $F(3) = 23 + 15 - 4 \cdot 3 = 26$ $\Rightarrow 26$. $\bigO{n}$ instead of $\bigO{n^2}$; values up to $10^5 \cdot 10^5 \cdot 100 = 10^{12}$ in intermediate products, so \texttt{long long} is the safe choice. Checked against computing all $n$ rotations directly.
\end{solutionbox}

\begin{solutionbox}[E3.5 Make Two Arrays Equal by Reversing Subarrays]
\textbf{The trap:} it looks like a search over sequences of reversals. But reversing a subarray of length 2 swaps two neighbours, and adjacent swaps can produce \emph{every} permutation (that is how bubble sort works). Reversals can do no more than permute. So the answer is \texttt{true} iff both arrays contain the same values with the same multiplicities.
\begin{center}
\begin{tikzpicture}[x=1cm, y=1cm]
    \dparray{x}{0.8}{}{2/acell, 4/aread, 1/aread, 3/acell}
    \dparray{y}{-0.4}{}{2/acell, 1/awrite, 4/awrite, 3/acell}
    \draw[dep, <->] (x1.north) to[bend left=40] (x2.north);
    \node[note, anchor=west, text width=6cm, align=left] at (4,0.2) {reverse $[1,2]$ = swap two neighbours; repeating this bubble-sorts any arrangement into any other with the same values};
\end{tikzpicture}
\end{center}
Values are at most 1000, so a counting array decides it in $\bigO{n + 1000}$ (Pattern~5): add 1 for each value of \texttt{target}, subtract 1 for each value of \texttt{arr}, and fail if a count becomes negative. Samples $\to$ true, true, false.
\end{solutionbox}

\begin{solutionbox}[E3.6 Sort the Array]
Let $l$ be the first index with $a_l > a_{l+1}$ (if none, the array is sorted: print \texttt{yes 1 1}). Everything before $l$ is already in place, so a valid segment must start at $l$. It must contain the whole decreasing run starting at $l$: extend $r$ while $a_r > a_{r+1}$. Reverse $[l, r]$ and check once whether the array is now increasing. \textbf{The trap} is to stop at the first ``bad pair'' or to reverse up to the \emph{last} bad pair without the final check: $[3,1,2,4]$ has its decreasing run $3,1$, reversing gives $[1,3,2,4]$, still unsorted $\Rightarrow$ \texttt{no}.
\begin{lstlisting}[style=cpp]
#include <bits/stdc++.h>
using namespace std;

int main() {
    int n;
    cin >> n;
    vector<int> a(n);
    for (auto& x : a) cin >> x;
    int l = 0;
    while (l + 1 < n && a[l] < a[l + 1]) l++;          // first break of the order
    if (l == n - 1) { cout << "yes\n1 1\n"; return 0; }
    int r = l;
    while (r + 1 < n && a[r] > a[r + 1]) r++;          // the whole decreasing run
    reverse(a.begin() + l, a.begin() + r + 1);
    for (int i = 0; i + 1 < n; i++)
        if (a[i] > a[i + 1]) { cout << "no\n"; return 0; }
    cout << "yes\n" << l + 1 << " " << r + 1 << "\n";
}
\end{lstlisting}
All four official samples reproduced ($[3,2,1] \to$ \texttt{1 3}, $[2,1,3,4] \to$ \texttt{1 2}, $[3,1,2,4] \to$ \texttt{no}, $[1,2] \to$ \texttt{1 1}); the yes/no answer matches trying every segment on random permutations with $n \le 6$. $\bigO{n}$ time.
\end{solutionbox}
